we have: @@M0@@.

\sgaStableTarget{sga6.e13.lem.2-11}\textbf{Lemma 2.11.} \emph{Let @@M1@@ be a noetherian scheme and @@M2@@ a projective @@M3@@-scheme, endowed with a very ample sheaf @@M4@@. We suppose that the fibers of @@M5@@ are geometrically integral, of the same dimension @@M6@@. Let @@M7@@ be a family of classes of invertible sheaves @@M8@@ on the fibers of @@M9@@. Then for @@M10@@ to be bounded, it is (necessary and) sufficient that, in the Hilbert polynomials @@M11@@ of the @@M12@@, all the coefficients @@M13@@ remain bounded. Moreover, if the fibers of @@M14@@ are also geometrically normal, then it suffices that the three coefficients @@M15@@ remain bounded.}

Indeed, let @@M16@@. If @@M17@@, one may replace @@M18@@ by the family of classes of the @@M19@@ by 2.3 and 2.10\,; one may thus suppose @@M20@@. Then @@M21@@ is a @@M22@@-family by 2.5, hence bounded if all the coefficients @@M23@@ remain bounded, by 1.13, and, when the fibers of @@M24@@ are geometrically normal, it suffices for this that the three coefficients @@M25@@ remain bounded, by 2.9.

\sgaStableTarget{sga6.e13.sec.3}\subsection*{3. General finiteness theorems}
\sgaStableTarget{sga6.e13.lem.3-1}\textbf{Lemma 3.1.} \emph{Let @@M26@@ be a scheme and @@M27@@ a quasi-compact and quasi-separated @@M28@@-scheme\,; let @@M29@@ be the Picard functor, (localized for the f.p.p.f. topology)}\SGAAuthorFootnote{*}{Cf. SGA 3 IV 6.3, and FGA no.\ 232 for the definition of @@M30@@, where, as indicated, the fpqc topology should here be replaced by the fppf topology.}.
\emph{Then for every filtered inverse system @@M31@@ of quasi-compact and quasi-separated @@M32@@-schemes
